\section{The binary coefficient identity}\label{sec:binary}

Proposition~\ref{prop:alternating} applies to every
\(x\in\Z[i][P]\) satisfying \(xx^*=u^2\). We now use the common array of signs to
relate three such elements. The following elementary map records the
first-order term of a unit \(1+tL\) and turns products into sums in
the residue field. Its common target will let us compare the two
divisibilities supplied by the sign array.

\begin{lemma}\label{lem:first-order}
Let \(O\) be a local domain with maximal ideal \(\n\), let
\(k=O/\n\), and let \(0\ne t\in\n\). Then \(1+tO\) is a subgroup of
\(O^\times\), and
\[
 \lambda_t:1+tO\longrightarrow(k,+),\qquad
 \lambda_t(1+tL)=L\bmod\n
\]
is a group homomorphism.
\end{lemma}

\begin{proof}
The expression \(1+tL\) is a unit because it has residue \(1\), and
\(L\) is unique because \(O\) is a domain and \(t\ne0\). The exact
identities
\[
 \frac{(1+tL)(1+tM)-1}{t}=L+M+tLM,\qquad
 \frac{(1+tL)^{-1}-1}{t}=-\frac{L}{1+tL}
\]
prove closure under multiplication and inversion. Reducing modulo
\(\n\) proves additivity and the inverse rule.
\end{proof}

\begin{proof}[Proof of Theorem~\ref{thm:main}]
The \(1\times1\) matrix \((1)\) is an example, as is the circulant
matrix with first row \((1,1,1,-1)\) at order four; its three
nontrivial autocorrelations vanish.

Suppose there were an example of another order. By
Proposition~\ref{prop:order}, its order is \(4u^2\) with \(u>1\) odd.
Write \(P=C_{u^2}\) and use the coprime decomposition
\(C_{4u^2}=C_4\times P\). For a generator \(Z\) of \(C_4\), write
the first-row element as
\[
 h=H_0+ZH_1+Z^2H_2+Z^3H_3,\qquad H_j\in\Z[P].
\]
This decomposition reindexes the coefficients of the row, so every
coefficient of every \(H_j\) is a sign. Turyn's Lemma~8
\cite[p.~337]{Turyn1965} studies the four coset counts under an
additional hypothesis on the ideals generated by the three nonprincipal
character values; here each \(H_j\) records the signs at
all elements of \(P\). The normalized partial evaluations at
\(Z=1,-1,i\) are
\begin{equation}\label{eq:three-elements}
 \begin{aligned}
 c&=\tfrac12h\vert_{Z=1}
   =\frac{H_0+H_1+H_2+H_3}{2},\\
 d&=\tfrac12h\vert_{Z=-1}
   =\frac{H_0-H_1+H_2-H_3}{2},\\
 g&=\tfrac12h\vert_{Z=i}
   =\frac{H_0-H_2}{2}+i\frac{H_1-H_3}{2}.
 \end{aligned}
\end{equation}
The fourth normalized evaluation, at \(Z=-i\), is the coefficientwise conjugate
of \(g\). The coefficients of \(c\) and \(d\) are integers because
their numerators are sums of four odd integers. The two half-differences
in \(g\) are integral as well. Hence \(c,d\in\Z[P]\) and
\(g\in\Z[i][P]\). Evaluation at \(Z=1,-1,i\) commutes with the involution,
so \(hh^*=4u^2\) gives
\begin{equation}\label{eq:three-norms}
 xx^*=u^2\qquad(x=c,d,g).
\end{equation}

Use the fixed notation \(I,B,\chi_S,x_S,\eps_S\) from
Section~\ref{sec:products}, and set
\[
 R=\frac{\Delta(d)}{\Delta(c)},\qquad
 W=\frac{\Delta(g)}{\Delta(c)}.
\]
Proposition~\ref{prop:alternating} shows that \(R\) and \(W\) are
roots of unity of odd order. Indeed, a quotient of roots whose orders
divide odd integers \(m\) and \(n\) has order dividing \(mn\).
This applies even when an individual alternating product is nontrivial.

Choose a maximal ideal \(\m\) of \(B\) containing \(2\), which exists
because \(B/2B\ne0\), and put
\[
 O=B_{\m},\qquad \n=\m B_{\m},\qquad k=O/\n.
\]
All residues in the rest of the proof lie in this one field \(k\).
Every \(c_S,d_S,g_S\) is a unit of \(O\), since its product with its
complex conjugate is the odd integer \(u^2\). The image of \(i\) in
\(k\) is \(1\): in characteristic two, \((i-1)^2=0\), and \(k\)
is a field. Thus \(2,1+i,1-i\) belong to \(\n\), and \(2\) and
\(1+i\) are nonzero in the characteristic-zero domain \(O\).

For every \(S\subseteq I\), define the local units
\begin{equation}\label{eq:ratios}
 r_S=\frac{d_S}{c_S},\qquad w_S=\frac{g_S}{c_S}.
\end{equation}
Expanding the alternating products gives
\begin{equation}\label{eq:ratio-products}
 R=\prod_{S\subseteq I}r_S^{\eps_S},\qquad
 W=\prod_{S\subseteq I}w_S^{\eps_S}.
\end{equation}
We now compare how these two families differ from \(1\).

Put \(b=(H_1+H_3)/2\in\Z[P]\). Direct subtraction in
\(\Z[i][P]\) gives the two exact differences
\begin{equation}\label{eq:raw-differences}
 d-c=-2b,\qquad
 g-c=(i-1)b-(H_2+iH_3).
\end{equation}
The common sign coefficients control the last bracket. Let
\(J=\sum_{z\in P}z\). Each coefficient of \(H_2+iH_3\) has the form
\(\alpha+i\beta\), with \(\alpha,\beta\in\{-1,1\}\), and differs
from \(1+i\) by an element of \(2\Z[i]\). Hence
\begin{equation}\label{eq:binary}
 H_2+iH_3=(1+i)J+2U,\qquad U\in\Z[i][P].
\end{equation}
The first difference in \eqref{eq:raw-differences} has a factor \(2\).
After substituting \eqref{eq:binary}, the second has a factor \(1+i\),
because \(i-1=i(1+i)\) and \(2=(1+i)(1-i)\).
Evaluating and dividing by the unit \(c_S\) therefore gives the
following integral first-order terms in \(O\):
\begin{align}
 L_{r,S}:=\frac{r_S-1}{2}
   &=-\frac{b_S}{c_S},\label{eq:lr}\\
 L_{w,S}:=\frac{w_S-1}{1+i}
   &=i\frac{b_S}{c_S}-\frac{J_S}{c_S}
     -(1-i)\frac{U_S}{c_S}.\label{eq:lw}
\end{align}
The right sides establish the integrality of the two quotients.
In particular
\[
 r_S\in1+2O,\qquad w_S\in1+(1+i)O.
\]
Adding the exact formulas makes the contribution of \(b_S\) disappear
after reduction:
\begin{equation}\label{eq:sum-first-order}
 L_{r,S}+L_{w,S}
 =-\frac{J_S}{c_S}-(1-i)\frac{b_S+U_S}{c_S}
 \equiv\frac{J_S}{c_S}\pmod\n.
\end{equation}

The displayed subgroup memberships give residue \(1\) for every
\(r_S\) and \(w_S\), and therefore for \(R\) and \(W\).
The orders of \(R\) and \(W\) are odd, so
Lemma~\ref{lem:torsion}(ii) in residue
characteristic two gives the exact equalities
\begin{equation}\label{eq:exact-products}
 R=W=1.
\end{equation}
A related character comparison of Leung and Schmidt also forces a root
of unity of odd order to equal one by reducing modulo two
\cite[author manuscript, proof of Theorem~3.5, p.~10]{LeungSchmidt2012}.
Lemma~\ref{lem:first-order} applies with \(t=2\) and with \(t=1+i\).
For the two families it gives
\[
 \lambda_2(r_S)=L_{r,S}\bmod\n,\qquad
 \lambda_{1+i}(w_S)=L_{w,S}\bmod\n.
\]
Both maps take values in \(k\), reducing their respective integral
quotients modulo the same maximal ideal \(\n\). Apply them to the
products in \eqref{eq:ratio-products}, using \eqref{eq:exact-products};
the homomorphism and inverse rules give
\[
 \sum_{S\subseteq I}\eps_S L_{r,S}\equiv0\pmod\n,\qquad
 \sum_{S\subseteq I}\eps_S L_{w,S}\equiv0\pmod\n.
\]
Adding these identities in \(k\) and using \eqref{eq:sum-first-order}
yields
\begin{equation}\label{eq:last-sum}
 \sum_{S\subseteq I}\eps_S\frac{J_S}{c_S}\equiv0\pmod\n.
\end{equation}

For \(S\ne\varnothing\), the character \(\chi_S\) is nontrivial.
Choose \(z\in P\) with \(\chi_S(z)\ne1\). Since multiplication by
\(z\) permutes the summands of \(J\), evaluation gives
\((\chi_S(z)-1)J_S=0\) in the domain \(B\), and hence \(J_S=0\).
For \(S=\varnothing\), one has \(J_{\varnothing}=|P|=u^2\) and
\(\eps_{\varnothing}=1\). The sum in \eqref{eq:last-sum} is therefore
the unit \(u^2/c_{\varnothing}\) of \(O\). Its residue cannot be zero.
This contradiction excludes every order greater than four.
\end{proof}

\begin{remark}
The hypothesis \(u>1\) is used when the nonempty prime set \(I\)
produces odd torsion in Proposition~\ref{prop:alternating}. At order
four, \(I\) is empty and there is still one subset,
\(S=\varnothing\). For the example \((1,1,1,-1)\), one has
\(c=d=1\) and \(g=i\), giving \(R=1\) and \(W=i\).
Although \(i\) has residue \(1\) in characteristic two, its order is
even, so the deduction \eqref{eq:exact-products} is unavailable.
\end{remark}

\section{Barker sequences}\label{sec:barker}

We finish by proving the length classification stated in
Corollary~\ref{cor:barker}. The new input is the circulant Hadamard
theorem for even lengths; the odd-length classification is classical.

\begin{proof}[Proof of Corollary~\ref{cor:barker}]
Suppose first that \(a\) is a Barker sequence of even length \(n>2\).
Its periodic autocorrelation at a nontrivial shift is
\[
 P_a(t)=\sum_{j=0}^{n-1}a_j a_{j+t\bmod n}
       =C_a(t)+C_a(n-t)\qquad(1\leq t<n).
\]
The product of the \(n\) sign summands in \(P_a(t)\) is
\((\prod_j a_j)^2=1\), so an even number of those summands are
negative. Consequently \(P_a(t)\equiv n\pmod4\).
Also \(C_a(t)\equiv n-t\pmod2\). For \(t=2\), both \(C_a(2)\)
and \(C_a(n-2)\) are even; the Barker bound makes them zero.
Thus \(P_a(2)=0\), and \(4\mid n\). For every nontrivial shift,
the Barker bounds give \(|P_a(t)|\leq2\), while the congruence gives
\(P_a(t)\equiv0\pmod4\). Hence all these periodic autocorrelations
vanish. The cyclic shifts of \(a\) form a real circulant Hadamard
matrix. This classical implication appears in Turyn and Storer
\cite[p.~395, footnote~2]{TurynStorer1961}; see also
\cite[p.~330]{Turyn1965}. Theorem~\ref{thm:main} now forces \(n=4\).
The only possible even lengths are therefore \(2\) and \(4\).

For odd \(n>1\), Turyn and Storer proved nonexistence above length
\(13\) \cite{TurynStorer1961}; the later proof of Schmidt and Willms
\cite[Theorem~1]{SchmidtWillms2016} gives the exact list
\(n\in\{3,5,7,11,13\}\). Examples exist at every listed length.
The following table uses the examples recorded in
\cite[Section~1]{SchmidtWillms2016}. Here \(+\) and \(-\) denote
\(1\) and \(-1\); direct substitution
in \(C_a(t)\) verifies the Barker property in each row.
\[
\begin{array}[b]{c|l@{\qquad}c|l}
 n & a & n & a\\ \hline
 2 & \texttt{++} & 7 & \texttt{+++{-}{-}+{-}}\\
 3 & \texttt{++-} & 11 & \texttt{+++{-}{-}{-}+{-}{-}+{-}}\\
 4 & \texttt{+++-} & 13 & \texttt{+++++{-}{-}++{-}+{-}+}\\
 5 & \texttt{+++-+} & &
\end{array}\qedhere
\]
\end{proof}

If length one is admitted, it is also possible: there are no nontrivial
shifts, so the Barker condition is vacuous.
